\section{El modelo es el producto, los datos son el modelo: Agentic Web y la línea principal de inversión}

La sección anterior señaló que el Online Learning requiere interacción y retroalimentación reales; esta sección lleva esa idea al producto, a los datos y a la estrategia de inversión. Guangmi (广密) insiste en que «el modelo es el producto, los datos son el modelo». La frase parece enrevesada, pero es muy concreta: la mayor mejora en la experiencia de los últimos años vino de la mejora de los modelos, y la raíz de la mejora de los modelos está, a su vez, en la mejora de la distribución de los datos. Quien consiga datos diferenciados podrá diferenciarse en el modelo y en el producto.

\begin{figure}[H]
\centering
\includegraphics[width=0.96\textwidth]{model-data-product-loop.png}
\caption{El modelo es el producto, los datos son el modelo: las tareas de los usuarios, los datos de interacción, la destilación de expertos, la mejora del modelo y la experiencia de producto forman un ciclo cerrado. Diagrama conceptual propio, elaborado a partir de la conversación de 00:43:05--00:56:45.}
\end{figure}

\begin{knowledgebox}{Lectura de la figura: la experiencia de producto vuelve a la distribución de los datos}
Las tareas de los usuarios exponen flujos de trabajo reales; los datos de interacción registran preferencias y puntos de falla; la destilación de expertos aporta procesos de alta calidad; la mejora del modelo amplía el límite de sus capacidades, y una mejor experiencia de producto atrae, a su vez, más tareas y más retroalimentación. Este ciclo cerrado es el sentido práctico de «el modelo es el producto, los datos son el modelo».
\end{knowledgebox}

\subsection{Un eje horizontal y uno vertical: destilación del conocimiento experto y Online Learning}

Guangmi resume el camino futuro como «un eje horizontal y uno vertical». El horizontal consiste en destilar el conocimiento de expertos humanos e incorporar más industrias, tareas y cadenas de herramientas a la distribución de datos del modelo; el vertical abarca el Online Learning, la Memory, los Agent proactivos y nuevas formas de interacción. El eje horizontal permite que la IA cubra más trabajo de oficina; el vertical podría permitir que la IA crezca de manera continua a partir de los usuarios y del entorno.

\begin{importantbox}{Un eje horizontal y uno vertical}
El eje horizontal responde a «qué industrias y tareas todavía no domina», y amplía el límite de capacidades mediante datos de expertos, flujos de trabajo reales y escenarios de cada industria; el eje vertical responde a «si puede volverse más inteligente cuanto más se usa», y se apoya en la Memory, el Online Learning, los Agent proactivos y la interacción a largo plazo.
\end{importantbox}

Esto tiene implicaciones directas para las empresas emergentes. Una UI delgada es fácil de copiar para las empresas de modelos base o las grandes plataformas; pero si se obtienen datos únicos, procesos de expertos, retroalimentación de usuarios y flujos de trabajo de una industria, es posible formar «activos de contenido» o «activos de proceso» sobre el modelo base. Por eso este episodio vuelve una y otra vez a dominios concretos como Coding, Office, Finance, salud, materiales y robótica.

\subsection{Agentic Web: ¿la verdadera Web3?}

Otro concepto importante del programa es la Agentic Web. Aquí Web3 no se refiere a la Web3 en el sentido de blockchain, sino a que, cuando el Agent se convierte en el nuevo punto de entrada, se reorganiza la relación entre las páginas web, las plataformas, los pagos, la búsqueda y las tareas de los usuarios. Antes, los puntos de entrada centrales de internet eran el cuadro de búsqueda, la App, el Feed y la Super App; el punto de entrada central de la Agentic Web podría ser «quiero lograr algo».

Si un Agent puede operar de forma continua, deja de ser solo una herramienta de consumo de contenido y se convierte en el intermediario que distribuye la atención, los pedidos, el tráfico y las transacciones. Para Uber, Ctrip (携程), la publicidad en buscadores, los fabricantes de teléfonos y las superaplicaciones, esta es una variable que obliga a renegociar quién controla el punto de entrada.

\begin{warningbox}{Las carencias de infraestructura de la Agentic Web}
La Agentic Web no se materializa solo con la capacidad del modelo. También necesita identidad, autorización, pagos, control de riesgos, verificación de humanos frente a máquinas, protocolos de sitios web, registros de auditoría y asignación de responsabilidades. Si estos componentes no avanzan al mismo ritmo, el Agent quedará atorado en el punto de «se puede demostrar, pero es difícil ejecutarlo a escala».
\end{warningbox}

\subsection{Estrategia de inversión: apostar donde el crecimiento tecnológico es más pronunciado}

Una vez enlazados el producto, los datos y la Agentic Web en un ciclo cerrado, esta sección pasa al marco de inversión. Aquí se habla de inversión no para dar recomendaciones de compra o venta, sino para usar la clasificación de un inversionista como vía inversa para entender la estructura de la industria: hacia qué capa fluye el dinero suele indicar qué capa es la más escasa y la que más amplifica el cambio tecnológico.

Guangmi resume la estrategia de inversión en IA así: invertir en los modelos más avanzados, invertir en el cómputo y la infraestructura de silicio que esos modelos necesitan, e invertir en los beneficios del derrame tecnológico de esos modelos. Este marco corresponde a tres clases de activos: empresas de modelos y de puntos de entrada como OpenAI/Google/Anthropic/ByteDance (字节); la oferta de base como Nvidia/TSMC/infraestructura de cómputo, y empresas que aprovechan el derrame tecnológico, como Cursor, Perplexity, Coding Agent, etiquetado de datos, Serving y herramientas multimodales. Aquí Perplexity es el nombre de una empresa de búsqueda y preguntas y respuestas con IA, no la perplexity/PPL de la evaluación de modelos de lenguaje; esta última es la métrica de «perplejidad», el exponencial de la entropía cruzada, que representa de forma aproximada la incertidumbre promedio del modelo ante el siguiente token.

El programa también menciona un cambio de cartera: de una apuesta más concentrada en OpenAI y ByteDance a una canasta formada por OpenAI, ByteDance, Google, Anthropic, Nvidia, TSMC y otras. La lógica detrás es la alternancia en el liderazgo y la incertidumbre sobre el paradigma: si es difícil prever al ganador final, conviene armar una cartera en torno a la línea principal en lugar de apostar por una sola empresa.

\begin{knowledgebox}{No es una recomendación de inversión, sino un mapa de la industria}
Estos apuntes conservan el marco de inversión del programa para entender la estructura de la industria: cómo se impulsan mutuamente los modelos, el cómputo, los puntos de entrada de las plataformas, el derrame hacia las aplicaciones y los servicios de datos. No es una recomendación sobre ningún valor ni sobre la valuación de ninguna empresa.
\end{knowledgebox}

\subsection{Los temas de 2026: Online Learning, trabajadores del conocimiento, multimodalidad y Agentic Web}

Esta sección reúne los juicios técnicos e industriales anteriores en una lista de observación para 2026. En lugar de preguntar «quién ganará al final», la pregunta más operativa es: qué temas mostrarán primero señales verificables y qué escenarios generarán primero ingresos y cambios en el comportamiento de los usuarios.

De cara a 2026, el programa enumera varios temas esperados. Primero, el Online Learning o señales del siguiente paradigma; segundo, Agent más proactivos y una Memory más sólida; tercero, capacidades L3/L4 parciales orientadas a los trabajadores del conocimiento, como Coding, Office/PPT/Excel, análisis financiero de inversiones y la obtención de información de cola larga; cuarto, un gran año para la multimodalidad; quinto, el impacto de la Agentic Web sobre el poder de distribuir el tráfico.

Aquí L3/L4 es una analogía tomada de los niveles de conducción autónoma: L3 puede entenderse como que el sistema asume la tarea principal en condiciones específicas, mientras la persona todavía debe supervisar o retomar el control; L4 significa completar la tarea con un alto grado de automatización dentro de escenarios acotados. Guangmi considera que el «L3/L4 parcial» es muy realista, mientras que el «L4 total» sigue siendo muy difícil. Dicho de otro modo, primero se obtiene valor económico en flujos de trabajo bien definidos, en lugar de esperar una AGI capaz de todo.

\subsection{Resumen de la sección}

«El modelo es el producto, los datos son el modelo» enlaza la tecnología, el producto y la inversión en un ciclo cerrado. La experiencia del modelo proviene de los datos, los datos provienen de tareas reales, las tareas reales provienen de los puntos de entrada del producto, y estos, a su vez, traen más retroalimentación. Lo decisivo en 2026 no es quién se declara más general, sino quién logra que los datos y el producto se refuercen de forma acumulativa en escenarios concretos.
